\section{Preliminaries}\label{Sect2}
The $\la$-function arises from the Kazhdan--Lusztig theory of Coxeter groups. We
review the relevant basic notions and facts in this section. Besides defining $\la$, we will recall the definitions of Kazhdan--Lusztig cells and the
``$\mu$-coefficients'' of Kazhdan--Lusztig polynomials. Both these notions will be key
to the proofs of our main theorems.

\subsection{Coxeter groups} 
\label{Coxeter def}
Throughout the article, $W$ shall denote a Coxeter
group with a finite generating set $S$ and Coxeter matrix $M=[m(s,t)]_{s,t\in S}$. Thus,
$m(s,s)=1$ for all $s\in S$, $m(s,t)=m(t,s)\in \Z_{\ge 2}\cup \{\infty\}$ for all
distinct $s,t\in S$, and $W$ is generated by $S$ subject to the relations
$(st)^{m(s,t)}=1$ for all $s,t$ for which $m(s,t)$ is finite. 
%
%
%
%


The defining data of each Coxeter group can be encoded via its \emph{Coxeter diagram}.
This is the weighted, undirected graph with vertex set $S$ and edge set
$\{\{s,t\}:s,t\in S,m(s,t)\ge 3\}$ such that each edge $\{s,t\}$ has weight
$m(s,t)$. Each edge is labelled by its weight except when the weight is $3$. 
A Coxeter group is called \emph{irreducible} if its Coxeter diagram is
connected; otherwise the group is \emph{reducible}. Note that any reducible
Coxeter group $W$ with Coxeter diagram $G$ is isomorphic to the direct product
of the Coxeter groups encoded by the connected components of $G$. 


Let $S^*$ be the free monoid generated by $S$. For any $w\in W$, we define the
\emph{length} of $w$, written $l(w)$, to be the minimum length of all words in
$S^*$ that express $w$. We call any such minimum-length word a \emph{reduced word} of
$w$. For any distinct $s,t\in S$, we call the relation
\[
 sts\cdots=tst\cdots 
\] 
where both sides have $m(s,t)$ factors a \emph{braid relation}. Since
$s^2=(ss)^{m(s,s)}=1$ for all $s\in S$, the
braid relation is equivalent to
the relation $(st)^{m(s,t)}=1$ from the definition of $W$. When $m(s,t)=2$, we call the relation $st=ts$ a \emph{commutation
relation}, for $s$ and $t$ commute. 

We can now recall the useful Matsumoto--Tits Theorem.
\begin{prop}[\cite{Tits};~{\cite[Theorem 1.9]{LG}}]
 \label{Matsumoto-Tits}
 Let $w\in W$. Then any pair of reduced words of $w$ can be obtained from each
 other by a finite sequence of braid relations.
\end{prop}


For more basic notions and facts about Coxeter groups such as 
the Bruhat order and its subword property, see~\cite{BB}.

\subsection{The \texorpdfstring{$\mathbf{a}$}{a}-function}
Let $W$ be an arbitrary Coxeter group. We recall Lusztig's definition of the function
$\la: W\ra \Z_{\ge 0}$ below.

Let $\cala=\Z[v,v\inverse]$. Following~\cite{LG}, we define the \emph{Hecke
algebra} of $W$ to be the unital $\cala$-algebra $H$ generated by the set
$\{T_s:s\in S\}$ subject to the relations 
\begin{equation}
 (T_s-v)(T_s+v\inverse)=0
 \label{eq:our normalization}
\end{equation}
for all $s\in S$ and 
\[
 T_sT_tT_s\cdots = T_tT_sT_t\cdots
\]
for all $s,t\in S$, where both sides have $m(s,t)$ factors.

It is well-known
that $H$ has a \emph{standard basis} $\{T_w:w\in W\}$ where
$T_w=T_{s_1}\cdots T_{s_q}$ for any reduced word $s_1\cdots s_q$ ($s_1,\ldots,
s_q\in S$) of $w$, as well as a \emph{Kazhdan--Lusztig basis} $\{C_w:w\in W\}$
with remarkable properties (see~\cite{EW}). Now let $h_{x,y,z} (x,y,z\in W)$ be the elements of $\cala$
such that
\[
 C_xC_y=\sum_{z\in W} h_{x,y,z}C_z
\]
for all $x,y$. By Lemma 13.5 of~\cite{LG}, for each $z\in W$, there exists a
unique integer $\la(z)\ge 0$ that satisfies the conditions
\begin{enumerate}
 \item $h_{x,y,z}\in v^{\la(z)}\Z[v\inverse]$ for all $x,y\in W$,
 \item $h_{x,y,z}\not\in v^{\la(z)-1}\Z[v\inverse]$ for some $x,y\in W$.
\end{enumerate}
This defines the function $\la: W\ra\Z_{\ge 0}$.


Elements of $\la$-value 0 or 1 are well understood in the following sense. 
\begin{prop}[{\cite[Proposition~13.7]{LG}}, {\cite[Corollary~4.10]{Xu}}]
 \label{subregular}
 Let $W$ be an arbitrary Coxeter group, and let $1_W$ be the identity of
 $W$. For all $w\in W$, we have
 \begin{enumerate}
 \item $\la(w)=0$ if and only if
 $w=1_W$.
 \item $\la(w)=1$ if and only if $w\neq 1_W$ and $w$ has a
 unique reduced word.
 \end{enumerate}
\end{prop}
\noindent Here, the set of non-identity elements with a unique reduced word is
known to be a \emph{two-sided Kazhdan--Lusztig
cell} (which we will define in the next subsection), and is sometimes called the \emph{subregular cell} (see~\cite{subregular}
and~\cite{Xu}). 
%
%
%

The following result classifies $\la(1)$-finite Coxeter groups, \ie Coxeter
groups with finite subregular cells, in
terms of Coxeter diagrams. We will use it in the proof of Theorem
1.3 in Section~6.
\begin{prop}[{\cite[Proposition~3.8]{subregular}}]
 \label{tree}
 Let $W$ be an irreducible Coxeter group with Coxeter diagram $G$. Then $W$ is
 $\la(1)$-finite if and
 only if $G$ is a tree and there is at most one edge of weight higher than 3
 in $G$.
\end{prop}
\noindent 

Besides the identity element and the elements in the subregular cell, it is
also easy to compute the $\la$-values
of products of commuting generators in a 
Coxeter group, thanks to the following two results of Lusztig.
\begin{prop}[{\cite[Section 14]{LG}}]
 \label{a local}
 Let $W$ be a Coxeter group with generating set $S$. Let $I\se S$ and let $W_I$ be the subgroup of $W$ generated by
 $I$. If $w\in W_I$, then $\la(w)$ computed in terms of $W_I$ is equal to
 $\la(w)$ computed in terms of $W$.
\end{prop}

\begin{rema}
 The above statement appears as part of Conjecture 14.2 in~\cite{LG}. However, it is
 known to hold in the setting of this paper, which is called the
 \emph{equal parameter} or the \emph{split} case (see~\cite{LG}, Section 15).
 The same remark applies to Proposition~\ref{constant}.
\end{rema}

\begin{prop}[{\cite[Proposition~13.8]{LG}}]
 \label{a for longest}
 Let $W$ be a finite Coxeter group, and let
 $w_0$ be the longest element of $W$. Then $\la(w_0)=l(w_0)$.
\end{prop}

\begin{coro}
 \label{product a}
 Let $W$ be a Coxeter group with generating set $S$. Let $I=\{s_1,s_2,\ldots,
 s_k\}$
 be a subset of $S$ such that $m(s_i,s_j)=2$ for all $1\le i<j\le k$, and let
 $w_0=s_1s_2\cdots s_k$. Then $\la(w_0)=k$.
\end{coro}
\begin{proof}
 The elements of $I$ commute with each other since $m(s_i,s_j)=2$
 for all distinct $i,j$, therefore the subgroup $W_I$ of $W$ generated by
 $I$ is isomorphic to the direct product of $k$ copies of the cyclic group of
 order 2. In particular, $W_I$ is finite. Furthermore, $w_0$ is clearly the
 longest element of $W_I$, therefore $\la(w_0)=l(w_0)=k$ by Propositions~\ref{a local} and~\ref{a for longest}.
\end{proof}


\subsection{Kazhdan--Lusztig cells}
We define the Kazhdan--Lusztig cells of a Coxeter group $W$ in this subsection. 
Let $H$ be the Hecke algebra of $W$, and let $\{C_w:w\in W\}$
be the Kazhdan--Lusztig basis of $H$. For each $x\in W$, let $D_x:H\ra \cala$
be the linear map such that 
\[
 D_{x}(C_y)=\delta_{x,y}
\]
for all $y\in W$, where $\delta$ is the Kronecker delta symbol. 
Furthermore, for $x,y\in W$,
\begin{enumerate}[leftmargin=2em]
 \item define $x\prec_L y$ if $D_{x}(C_sC_y)\neq 0$ for some
 $s\in S$;
 \item define $x\le_L y$ if there is a sequence $x=z_1,z_2,\ldots,z_n=y$ in
 $W$ such that $z_{i}\prec_L z_{i+1}$ for all $1\le i\le n-1$;
 \item define $x\sim_L y$ if $x\le_L y$ and $y\le_L x$.
\end{enumerate}
By the construction, $\sim_L$ defines an equivalence
relation on $W$. 
%
%
We call the equivalence classes the
\emph{left Kazhdan--Lusztig cells}, or simply the \emph{left cells}, of $W$,
and we define the \emph{right (Kazhdan--Lusztig) cells}
and $\emph{two-sided (Kazhdan--Lusztig) cells}$ of $W$ similarly. Here, to
define the two-sided cells, start by declaring $x\prec_{LR} y$ if either
$D_{x}(C_sC_y)\neq 0$ for some $s\in S$ or $D_x(C_yC_s)\neq 0$ for some $s$
(\ie if either $x\prec_L y$ or $x\prec_R y$).
Note that each two-sided cell of $W$ must be a union of left cells as well as a
union of right cells. 

The Kazhdan--Lusztig cells of $W$ have the following key connection with the
$\la$-function on $W$.

\begin{prop}[{\cite[Section 14]{LG}}]
 \label{constant}
 Let $x,y\in W$. If $x\le_{LR} y$, then $\la(x)\ge \la(y)$. In particular, if
 $x\sim_{LR} y$, then $\la(x)=\la(y)$.
\end{prop}

By the proposition, one way to establish that an element $x\in W$ has a certain
$\la$-value is to find another element $y$ of that $\la$-value and prove that
$x$ and $y$ are in the same cell. We will repeatedly use this strategy in
Sections~5.2 and~5.3.

As we may see from their construction, the key to understanding Kazhdan--Lusztig cells lies in
understanding the products of the form $C_sC_y$. These products are controlled
by the \emph{Kazhdan--Lusztig polynomials}, which are defined to be the elements $p_{x,y}\in \cala\, (x,y\in W)$
such that 
\[
 C_y=\sum_{x\in W} p_{x,y}T_x
\]
for all $y\in W$, where the elements $\{T_w:w\in W\}$ form the standard basis
of $T$. More precisely, for each $x,y\in W$, let $\mu_{x,y}$ be the coefficient of the term
$v^{-1}$ in $p_{x,y}$, then we have the following formulae.
\begin{prop}[{\cite[Theorem 6.6]{LG}}]
 \label{KL basis mult}
 Let $y\in W$, $s\in S$, and let $\le$ be the Bruhat order on $W$. Then in the
 Hecke algebra $H$ of $W$,
 \begin{align*}
 C_s C_y&= 
 \begin{cases}
 (v+v\inverse) C_y&\quad\qquad\text{if}\quad sy<y,\\
 C_{sy}+\sum\limits_{x:sx<x<y} \mu_{x,y}C_x&\quad\qquad \text{if} \quad
 sy>y,
 \end{cases}\\
 C_y C_s&= 
 \begin{cases}
 (v+v\inverse) C_y&\quad \text{if}\quad ys<y,\\
 C_{ys}+\sum\limits_{x:xs<x<y} \mu_{x^{-1},y^{-1}}C_x&\quad \text{if}
 \quad ys>y.
 \end{cases}
 \end{align*}
\end{prop}

\begin{rema}
 \label{symmetry}
 It is known that $\mu_{x,y}=\mu_{x\inverse,y\inverse}$ for any $x,y\in W$
 (see Section 5.6 and Corollary 6.5 of~\cite{LG}),
 therefore the last formula in the proposition also holds with $\mu_{x,y}$ in
 place of $\mu_{x\inverse,y\inverse}$.
\end{rema}

\begin{rema}
 \label{conversion}
 The paper~\cite{KL} uses a normalization of the Hecke algebra that is
 different from ours, namely, it uses the relation $(T_s+1)(T_s-q)=0$ in place
 of our Equation~\eqref{eq:our normalization}. Consequently, the
 Kazhdan--Lusztig polynomials $P_{x,y}$ obtained in~\cite{KL}, which are
 polynomials in $q$, do not exactly
 agree with our Kazhdan--Lusztig polynomial $p_{x,y}$. However, it is
 straightforward to 
 check that we may convert $P_{x,y}$ to $p_{x,y}$ by first substituting
 $q$ by $v^2$ in $P_{x,y}$ and then multiplying the result by $v^{l(x)-l(y)}$. 
 In particular, our definition of the numbers $\mu_{x,y}$ agrees with that in
~\cite{KL}.
\end{rema}

\noindent The $\mu$-coefficients are often called the ``leading coefficients
of Kazhdan--Lusztig polynomials'' in the literature. Note that the elements $C_s (s\in S)$ generate $H$ by Proposition~\ref{KL basis mult}, so in a sense the $\mu$-coefficients control the
multiplication of the Kazhdan--Lusztig basis elements in the Hecke algebra. As
such, they also lead to an alternative characterization of the relations
$\prec_L$ and $\prec_R$: for each $y\in W$, define the \emph{left descent set} and
\emph{right descent set} of $y$ to be the sets
\[
 \mathcal{L}(y)=\{s\in S:sy<y\},\qquad \mathcal{R}(y)=\{s\in S:ys<y\},
\]
respectively. Then the following proposition holds.

\begin{prop}
 \label{alternative prec}
 Let $x,y\in W$. Then
 \begin{enumerate}
 \item\label{prop2.12_1} $x\prec_L y$ if and only if one of the following conditions holds:
 \textup{(}a\textup{)} $x=y\neq 1_W$;
 \textup{(}b\textup{)} $x=sy$ for some $s\notin
 \mathcal{L}(y)$; \textup{(}c\textup{)} $x<y$, $\mathcal{L}(x)\not\se \mathcal{L}(y)$, and
 $\mu_{x,y}\neq 0$.
 \item\label{prop2.12_2} $x\prec_R y$ if and only if one of the following conditions holds:
 \textup{(}a\textup{)} $x=y\neq 1_W$;
 \textup{(}b\textup{)} $x=ys$ for some $s\notin
 \mathcal{R}(y)$; \textup{(}c\textup{)} $x<y$, $\mathcal{R}(x)\not\se \mathcal{R}(y)$, and
 $\mu_{x,y}\neq 0$.
 \end{enumerate}
\end{prop}
\begin{proof}
 By Proposition~\ref{KL basis mult}, we have $D_x(C_sC_y)\neq 0$ for some $s\in S$ if and only if one of the following occurs:
 \begin{enumerate}[(\alph*)]
 \item %
 $x=y\neq 1_W$, so that $\mathcal{L}(y)\neq \emptyset$ and $D_x(C_sC_y)\neq 0$ for each $s\in
 \mathcal{L}(y)$;
 \item %
 $C_x$ appears in $C_sC_y$ for some $s\not\in \mathcal{L}(y)$, with $x=sy$;
 \item %
 $C_x$ appears in $C_sC_y$ for some $s\not\in \mathcal{L}(y)$, and $x$
 satisfies $x<y, sx<x$ and $\mu_{x,y}\neq 0$. Note that in this case
 we have $\mathcal{L}(x)\not\se\mathcal{L}(y)$ as $s\in
 \mathcal{L}(x)\setminus \mathcal{L}(y)$. Conversely, 
 if $x<y$, $\mathcal{L}(x)\not\se\mathcal{L}(y)$, and $\mu_{x,y}\neq 0$,
 then 
 $D_x(D_sD_y)\neq 0$ for each $s\in \mathcal{L}(x)\setminus
 \mathcal{L}(y)$ by Proposition~\ref{KL basis mult}.
 \end{enumerate}
 Statement~\eqref{prop2.12_1} now follows. The proof of~\eqref{prop2.12_2}
 is similar.
\end{proof}

In Propositions~\ref{star and mu} and~\ref{star relation}, we will describe
ways to compute certain
$\mu$-coefficients combinatorially without referring to the Hecke algebra.
This will allow us to avoid difficult computations of Kazhdan--Lusztig
polynomials and understand Kazhdan--Lusztig cells by using only the combinatorics of
Coxeter groups.

To end this section, we record several facts for future use.

\begin{coro} \label{weak order} Let $x,y\in W$, and let $s\in
S$. 
\begin{enumerate} 
\item\label{coro2.13_1} We have $sy\le_L y$ if $sy>y$, and $ys\le_R y$ if $ys>y$;
\item\label{coro2.13_2} If there exist elements $u,v\in W$ such that $y=uxv$ and
 $l(y)=l(u)+l(x)+l(v)$ where $l$ is the length function on
 $W$ \textup{(}see Section~\ref{Coxeter def} for the definition of
 $l$\textup{)}, then
 we have $y\le_{LR}x$ and $\la(y)\ge \la(x)$. 
 \end{enumerate} 
 \end{coro}
 \begin{proof} 
 This is a simple corollary of Propositions~\ref{KL basis
 mult} and~\ref{constant}. Note that~\eqref{coro2.13_2} follows from repeated application of~\eqref{coro2.13_1} and
 Proposition~\ref{constant}, hence it suffices to prove~\eqref{coro2.13_1}. Suppose
 $sy>y$. Then $D_{sy}(C_sC_y)=1$ by Proposition~\ref{KL basis mult},
 therefore $sy\prec y$ and $sy\le_L y$ by definition. Similarly,
 we have $ys\le_R y$ if $ys>y$. \end{proof}


 \begin{prop}[{\cite[Proposition~5.4]{LG}}]
 \label{degree bound}
 Let $x,y\in W$. If $x\le y$, then $p_{x,y}=v^{-l(y)+l(x)}\mod
 v^{-l(y)+l(x)+1}\Z[v]$.
 \end{prop}

 \begin{coro}
 \label{difference 1}
 Let $x,y\in W$. If $x\le y$ and $l(x)=l(y)-1$, then $p_{x,y}=v\inverse$
 and hence $\mu_{x,y}=1$.
 \end{coro}
 \begin{proof}
 This is immediate from the
 well-known fact that $p_{x,y}\in \Z[v\inverse]$ (see \cite[Section~5.3]{LG}) and Proposition~\ref{degree bound}.
 \end{proof}
 %
 %
 %
 %
 %
 %
 %
 %
 %
 %
 %
 %
 %

 \begin{prop}[{\cite[Fact 5]{Warrington}}]
 \label{extremal}
 Let $x,y\in W$ be such that $l(x)<l(y)-1$. If $\mathcal{L}(y)\not\se
 \mathcal{L}(x)$ or $\mathcal{R}(y)\not\se\mathcal{R}(x)$, then $\mu(x,y)=0$.
 \end{prop}

 %
 %
 %
 %
 %
 %
 %
 %
 %
 %
 %
 %
 %

 \section{Tools for computation of \texorpdfstring{$\mathbf{a}$}{a}}
 \label{sec:tools}
 We introduce our main tools for verification and computation of
 $\la$-values in this section. The first tool is the so-called
 generalized star operations, which we will often use to show two elements
 are in a same Kazhdan--Lusztig cell and hence of the same $\la$-value.
 The second tool involves heaps of fully commutative elements and will
 allow us to directly compute $\la$-values in certain cases.
 

 \subsection{Generalized star operations}
 \label{star}
 We review the notion of a
 \emph{generalized star operation} in this subsection. We highlight a
 direct connection between the operation and Kazhdan--Lusztig cells, then describe
 a more subtle recurrence relation involving the operation and the
 $\mu$-coefficients from Proposition~\ref{KL basis mult}. 

 Let $W$ be an arbitrary Coxeter group, and let $s,t\in S$ be a pair of
 generators of $W$ with $3\le m(s,t)<\infty$. Set $I=\{s,t\}$, let $W_I=\langle
 s,t\rangle$, the subgroup of $W$ generated by $s$ and $t$, and set ${}^I W=\{w\in
 W: \mathcal{L}(w)\cap I=\emptyset\}$. It is known that every $w\in W$ admits a unique
 factorization $w=w_I\cdot {}^I w$, called a \emph{coset decomposition}, where
 ${}^I w\in {}^I W$ and $w_I\in W_I$; moreover, we have $l(w)=l(w_I) +
 l({}^Iw)$ in this case. (For proofs
 of these facts and an algorithm to compute the factors $w_I$ and
 ${}^I w$, see~\cite[Proposition~2.4.4]{BB}.) Consider the following situations:
 \begin{enumerate}
 \item\label{sec3.1_1} $w_I=1$;
 \item\label{sec3.1_2} $w_I$ is the longest element $sts\cdots$ of length $m(s,t)$ in
 $W_I$;
 \item\label{sec3.1_3} $w$ is one of the $(m-1)$ elements $s\cdot {}^Iw, ts\cdot {}^Iw,
 sts\cdot {}^I
 w,tsts\cdot {}^I w,\ldots$;
 \item\label{sec3.1_4} $w$ is one of the $(m-1)$ elements $t\cdot {}^Iw, st\cdot {}^Iw,
 tst\cdot {}^I 
 w,stst\cdot {}^I w,\ldots$.
 \end{enumerate}
 We call the sequences appearing in~\eqref{sec3.1_3} and~\eqref{sec3.1_4} \emph{left
 $\{s,t\}$-strings} or $\emph{left $I$-strings}$, or simply
 \emph{left strings} if the pair $\{s,t\}$ is clear from context. For any element $w$ in a left
 $\{s,t\}$-string other than the longest, we define ${}^* w$ to be the
 element to the right of $w$. Otherwise, we leave ${}^* w$ undefined. We call
 the map $w\mapsto {}^* w$ the \emph{upper left star operation with respect
 to $I$}. 

 Similarly, we define the \emph{lower left
 star operation} to be the operation $w\mapsto {}_* w$ where $w$ is an element
 in a left string other than the shortest and ${}_* w$ is the element to the
 left of $w$ in the same string. In addition, we say $w$ is \emph{left star reducible to ${}_*
 w$ with respect to $I$} whenever the latter is defined. More generally,
 dropping the reference to a particular pair of generators, we
 say $y$ is \emph{left star reducible} to $x$ for $x,y\in W$ if there is a sequence
 $x=z_1,z_2,\ldots,z_n=y$ in $W$ such that for each $1\le i\le n-1$, there is
 some pair $I_i=\{s_i,t_i\}\se S$ with $3\le m(s_i,t_i)<\infty$ such that
 $z_{i+1}$ is left star reducible to $z_i$ with respect to $I_i$.



 The concepts and notations above have obvious right-handed counterparts,
 where the coset decompositions to be considered are of the form
 $w=w^I\cdot w_I$ where $w_I\in W_I$ and the factor $w^I$ is from the set $W^I:=\{w\in
 W:\mathcal{R}(w)\cap I=\emptyset\}$. We
 refer to the two types of left star operations and their right-handed
 counterparts collectively as
 \emph{generalized star operations}. Finally, for $x,y\in W$, we say that $y$
 is \emph{star reducible} to $x$ if there is a sequence
 $x=z_1,z_2,\ldots,z_n=y$ in $W$ such that $z_{i+1}$ is either left reducible
 or right reducible to $z_i$ for each $1\le i\le n-1$.

 \begin{rema}
 \label{original star}
 For each pair $I=\{s,t\}\se S$ with $m(s,t)=3$ and each member of a left $\{s,t\}$-string, only one of the lower and upper left star
 operations is defined for
 each member of a left $\{s,t\}$-string. The one that does is simply called
 the \emph{left star operation} in the paper~\cite{KL} where the operation
 was first introduced by Kazhdan and Lusztig. Similarly, it makes sense to
 simply speak of a \emph{right star operation} with respect to $I$. 
 \end{rema}


 \begin{exam}
 \label{star example}
 Let $W$ be a Coxeter group with generating set $S=\{a,b,c\}$. Suppose
 $m(a,b)=3$, $m(b,c)=4$, $m(a,c)=2$, let $I=\{a,b\}, J=\{b,c\}$, and let $x=abcab$. Then with respect to $I$, the coset decompositions of
 $x$ are given by
 \[x= x_I\cdot {}^I x= aba\cdot cb, \quad x= x^I \cdot x_I = abc\cdot ab.\]
 It follows that $x$ is not in a left $I$-string but is in a right $I$-string.
 Moreover, as pointed out in Remark~\ref{original star}, only one right star operation with respect
 to $I$ is defined on $x$ since $m(a,b)=3$: only the lower star
 operation is defined, and we have $x_*=abca$. With respect to
 $J$, the coset decompositions of $x$ are given by 
 \[
 x=x_J\cdot {}^J x=b\cdot abcb, \quad x=x^J\cdot x_J=ba\cdot bcb.\]
 It follows that $x$ is both in a left $J$-string and in a right $J$-string.
 Moreover, we have ${}^{*}x=cbabcb$ and $x_*=babc$, but ${}_*x$ and
 $x^*$ are not defined.
 \end{exam}


 Generalized star operations are intimately related to Kazhdan--Lusztig
 cells:

 \begin{prop} 
 \label{star and cell} 
 Let $W$ be an arbitrary Coxeter group, and let $I=\{s,t\}$ be a pair of
 generators of $W$ for which $3\le m(s,t)<\infty$. Then the following hold,
 where all star operations are performed with respect to $I$.
 \begin{enumerate} 
 \item\label{prop3.3_1} Let $y$ be an element of a left $\{s,t\}$-string such that 
 ${}_*y$ makes sense, then $y\sim_{L}
 {}_* y$. 
 \item\label{prop3.3_2} Let $y$ be an element of a right $\{s,t\}$-string such that 
 $y_*$ makes sense, then $y\sim_{R} y_*$. 
 \end{enumerate} 
 \end{prop}

 \begin{rema}
 The above facts are well-known to experts, but we have
 not found a reference stating it explicitly in this way, so we include a
 brief proof below. 
 \end{rema}
 \begin{proof}
 We first prove~\eqref{prop3.3_1}. Without loss of generality, suppose $I\cap
 \mathcal{L}(y)=\{s\}$. Then the definition of left strings guarantees that
 $I\cap\mathcal{L}({}_*y)=\{t\}$. Since ${}_*y<y$, $t\in \mathcal{L}({}_*
 y)\setminus\mathcal{L}(y)$ and $\mu_{{}_*y,y}=1$ by Corollary
~\ref{difference 1}, we have ${}_* y\le_L y$ by Proposition~\ref{alternative
 prec}. On the other hand, $y\le_L {}_* y$ by Corollary~\ref{weak order},
 therefore $y\sim_L {}_* y$. The proof of~\eqref{prop3.3_2} is similar.
 \end{proof} 

 \begin{coro}
 \label{star a}
 Let $x,y\in W$. If $y$ is star reducible to $x$, then $\la(x)=\la(y)$.
 \end{coro}
 \begin{proof}
 Suppose $y$ is star reducible to $x$. Then $x\sim_{LR} y$ by repeated
 application of Proposition~\ref{star and cell}, therefore
 $\la(x)=\la(y)$ by Proposition~\ref{constant}.
 \end{proof}

 Generalized star operations are also connected with $\mu$-coefficients:

 \begin{prop}[{\cite[Theorem 4.2]{KL}}]
 \label{star and mu}
 Let $W$ be an arbitrary Coxeter group, and let $I=\{s,t\}$ be a pair of
 generators of $W$ for which $m(s,t)=3$. Then the following hold, where all
 star operations are performed with respect to $I$.
 \begin{enumerate}
 \item Let $x,y\in W$ be elements of left $\{s,t\}$-strings such that
 $xy^{-1}\notin W_I$. Then $\mu(x,y)=\mu(*x,*y)$, where $*\alpha$ stands for the
 result of applying the left star operation on $\alpha$ for each string
 $\alpha$ \textup{(}see Remark~\ref{original star}\textup{)};
 \item Let $x,y\in W$ be elements of right $\{s,t\}$-strings such that
 $x^{-1}y\notin W_I$. Then $\mu(x,y)=\mu(x*,y*)$, where $\alpha*$ stands for the
 result of applying the right star operation on $\alpha$ for each string
 $\alpha$ \textup{(}see Remark~\ref{original star}\textup{)}.
 \end{enumerate}

 \end{prop}

 \begin{prop}[{\cite[Section 10.4]{Mu}}; {\cite[Proposition~5.9]{Green_Jones}}]
 \label{star relation}
 Let $W$ be an arbitrary Coxeter group, and let $I=\{s,t\}$ be a pair of
 generators of $W$ for which $3\le m(s,t)<\infty$. Then the following
 hold,
 where all star operations are performed with respect to~$I$.

 \begin{enumerate}
 \item\label{prop3.7_1} 
 Let $x,y\in W$ be
 elements of left $\{s,t\}$-strings such that $\mathcal{L}(x)\cap I\neq
 \mathcal{L}(y)\cap I$. Then
 \[
 \mu({}_* x,y)+\mu({}^* x,y)=\mu(x,{}_* y)+\mu(x, {}^* y);
 \]

 \item\label{prop3.7_2} 
 Let $x,y\in W$ be
 elements of right $\{s,t\}$-strings such that $\mathcal{R}(x)\cap I\neq
 \mathcal{R}(y)\cap I$. Then
 \[
 \mu(x_*,y)+\mu(x^*,y)=\mu(x,y_*)+\mu(x, y^*).
 \]
 \end{enumerate}
 Here, we define $\mu(\alpha,\beta)=0$ if either $\alpha$ or $\beta$ is an undefined symbol.
 \end{prop}


 In Section~5.2, we will frequently use the two propositions above to
 compute certain $\mu$-coefficients $\mu_{x,y}$ recursively, then use Proposition~\ref{alternative prec} to conclude that $x\sim_{L}y$ or $x\sim_R y$. This provides
 a very useful connection, albeit a less direct one than Proposition~\ref{star
 and cell}, between generalized star operations and cells. 

 \subsection{Full commutativity and heaps} 
 \label{fc}
 Let $W$ be an arbitrary Coxeter group. In this subsection,
 we show that any element with $\la$-value 2 must be
 \emph{fully commutative} in the sense of~\cite{FC}. We then recall a
 combinatorial characterization of the $\la$-values of fully commutative
 elements in a Weyl or affine Weyl group in terms of \emph{heaps}. This
 characterization will allow us to compute certain $\la$-values without
 recourse to Kazhdan--Lusztig theory in Section~5.1.

 An element $w\in W$ is said to be \emph{fully commutative} if any pair of
 reduced words of $w$ can be obtained from each other by means of only
 commutation relations. It is well-known that $w$ is fully commutative if
 and only if no reduced word of $w$ contains a
 contiguous subword of $sts\cdots$ of length $m(s,t)$ where $s,t\in S$
 and $m(s,t)\ge 3$ (see~\cite[Proposition~2.1]{FC}).

 \begin{rema}
 \label{removal}
 Let $w$ be a fully commutative element with a reduced word $w=stw'$ where
 $l(w)=l(w')+2$ and $m(s,t)\ge 3$. Consider the coset decomposition
 $w=w_I\cdot {}^I w$ with respect to the pair $I=\{s,t\}$. Since $w$ is
 fully commutative, $w_I$ cannot be the word $sts\cdots$ of length
 $m(s,t)$, therefore $w$ is an element of a left $\{s,t\}$-string, with ${}_* w=tw'$ with respect to $I$. That is, whenever a reduced word
 of a fully commutative element starts with a pair of letters $s,t\in S$
 with $m(s,t)\ge 3$, the lower left star operation with respect to
 $\{s,t\}$ simply removes the leftmost letter of $w$. Similarly,
 whenever a reduced word of a fully commutative element ends with a pair
 of letters $s,t\in S$ with $m(s,t)\ge 3$, the lower right star
 operation with respect to $\{s,t\}$ simply removes the rightmost letter
 of $w$. \end{rema}


 Problems related to fully commutative elements, such as the classification of
 Coxeter groups with finitely many fully commutative elements, the
 enumeration of fully commutative elements for those groups, and 
 connections fully commutative elements have to Kazhdan--Lusztig cells and
 so-called \emph{generalized
 Temperley--Lieb algebras}, have been studied
 extensively; see, for example,~\cite{BJN}, \cite{Fan}, \cite{FCcells}, \cite{Shi}, \cite{FC} and~\cite{FC2}.
 Fully commutative elements also provide a suitable framework for studying elements of $\la$-value 2 because of the following fact.

 \begin{prop}
 \label{a2 is fc}
 Let $w\in W$. If $\la(w)=2$, then $w$ is fully commutative. 
 \end{prop}
 \begin{proof}
 We prove the contrapositive of the statement, \ie that if $w$ is not fully
 commutative, then $\la(w)\neq 2$. 

 Suppose $w$ is not fully commutative. Then $w$ can be written in the form
 $w=uxv$ where $l(w)=l(u)+l(x)+l(v)$ and $x$ is of the form
 $x=sts\cdots$ with $s,t\in S, m(s,t)\ge 3$ and $l(x)=m(s,t)$. By Propositions~\ref{a local} and $\ref{a for longest}$, we have $\la(x)=m(s,t)$ in this case, therefore $\la(w)\ge \la(x)=m(s,t)\ge 3$ by
 Corollary~\ref{weak order}. This completes the proof.
 \end{proof}

 Next, we define the \emph{heap} of an arbitrary word $s_1s_2\cdots s_q$ in the free
 monoid $S^*$: this is the poset $([q],\preccurlyeq)$ where
 $[q]=\{1,2,\ldots,q\}$ and $\preccurlyeq$ is the partial order on
 $[q]=\{1,2,\ldots,q\}$ obtained via the reflexive transitive closure of the relations
 \[
 i\prec j \quad\text{if}\quad i<j\quad\text{and}\quad m(s_i,s_j)\neq 2.
 \]
 In particular, $i\prec j$ if $i<j$ and $s_i=s_j$. We refer to the
 generator $s_i$ as the \emph{label} of $i$ for each $i\in [q]$. It is well-known that the heaps of
 any two words in $S^*$ related by a commutation relation from $W$ are
 isomorphic as posets (see~\cite{FC}, Section 2.2), therefore for any fully commutative element $w\in W$,
 it makes sense to define the \emph{heap of $w$} to be the heap of any reduced
 word of $w$. In this case, we denote the heap of $w$ by $H(w)$. On the other
 hand, given any word in $S^*$, there is also a criterion for
 determining if its heap is that of a fully commutative element:

 \begin{prop}[{\cite[Proposition~3.3]{FC}}]
 \label{fc criterion}
 The heap $P$ of a word $s_1s_2\cdots s_q$ in $S^*$ is the heap of some fully commutative
 element in $W$ if and only if 
 \begin{enumerate}
 \item There is no covering relation $i\prec j$ such that $s_i=s_j$;
 \item There is no convex chain $i_1<i_2\cdots <i_m$ in $P$ such that
 $s_{i_1}=s_{i_3}=\cdots = s$ and $s_{i_2}=s_{i_4}=\cdots=\cdots=t$,
 where $s,t\in S$ and $m = m(s,t)\ge 3$.
 \end{enumerate}
 \end{prop}

 There is an intuitive way to visualize heaps of words in $S^*$. Consider the lattice $S\times \N$, with $S$
 indexing the \emph{columns} of the lattice and $\N$ indexing the \emph{levels}
 or
 \emph{heights}. We say two
 columns $s,t$ are \emph{adjacent} if the corresponding vertices are adjacent in the
 Coxeter graph, \ie if $m(s,t)\ge 3$.

 For any word $s_1s_2\cdots s_q\in S^*$, we may
 embed its heap $P$ as a set of lattice points in $S\times \N$ as follows:
 read the word from left to right, and drop a point in the column representing
 $s_i$ as we read each letter. Here, we envision each point as being under the
 influence of ``gravity'' in the sense that the point must fall to the lowest
 possible row in its column subject to one condition, namely, it must fall
 higher than every point that was placed before it in the same column
 or in an adjacent column. We define the index of this lowest possible row
 to be the \emph{level} of the point. 
 
 \begin{rema}
 A poset $P$ is said to be \emph{ranked} if there exists a function
 $\rho: P\ra \Z$, called a \emph{rank function for $P$}, such that
 $\rho(b)=\rho(a)+1$ whenever $a,b$ are elements in $P$ for which $a\prec b$
 is a covering relation in
 $P$. It is worth noting that for a fully commutative element
 $w$ in
 $W$, while we have described how to assign each element in the heap $P$
 of $w$ a
 well-defined level, the poset $P$ is not
 necessarily ranked, \ie the level function may not be a rank function.
 For an example where $P$ is not ranked and for criteria for $P$ to be
 ranked, see~\cite{Green_interval}.
 \end{rema}

 
 As we embed $P$ in the lattice $S\times \N$, after a point $k\in P$ ($1\le k\le q$) falls into position, it is customary to
 label the point with $s_k$ rather than $k$. Furthermore, to indicate the covering relations of the heap, we connect the point with
 edges to the highest existing points in its column and its adjacent
 columns. The resulting graph therefore recovers the Hasse diagram
 of the poset $P$. 
 For example, in Figure~\ref{fig:example}, the picture on the right shows the heap of
 the element $abcabd$ in the Coxeter group whose Coxeter diagram is drawn on the
 left. Note that
 all reduced words of a fully commutative element result in an identical
 graph when we embed them in $S\times \N$, so
 we may identify the element with its embedding. For more on the lattice
 embeddings of heaps, see~\cite{BJ}.

 

 \begin{figure}[htb]\centering
 \centering
 \begin{tikzpicture}
 \node[main node] (a) {};
 \node[main node] (b) [right=1cm of a] {};
 \node[main node] (c) [right=1cm of b] {};
 \node[main node] (d) [right=1cm of c] {};

 \node (aa) [below=0.2cm of a] {$a$};
 \node (bb) [below=0.1cm of b] {$b$};
 \node (cc) [below=0.2cm of c] {$c$};
 \node (dd) [below=0.1cm of d] {$d$};

 \path[draw]
 (b)--(c)--(d)
 (a) edge node [above] {$4$} (b);

 \node[main node] (1) [below right=2cm and 5cm of a] {};
 %
 \node[main node] (3) [above right=1cm and 1cm of 1] {};
 \node[main node] (4) [above left=1cm and 1cm of 3] {};
 \node[main node] (5) [above right=1cm and 1cm of 3] {};
 \node[main node] (6) [above left=1cm and 1cm of 5] {};
 \node[main node] (7) [above right=1cm and 1cm of 5] {};

 \node (11) [above=0.1cm of 1] {$a$};
 %
 \node (33) [above=0.1cm of 3] {$b$};
 \node (44) [above=0.1cm of 4] {$a$};
 \node (55) [above=0.1cm of 5] {$c$};
 \node (66) [above=0.1cm of 6] {$b$};
 \node (77) [above=0.1cm of 7] {$d$};

 \path[draw]
 (1)--(3)--(4)--(6)--(5)--(7)
 (3)--(5);

 \end{tikzpicture}

 \caption{Lattice embedding of a heap.}
 \label{fig:example}
 \end{figure}
 Note that the criterion from
 Proposition~\ref{fc criterion} can now be translated as follows.
 \begin{prop}
 \label{fc heap criterion}
 The heap $P$ of a word in $S^*$ is the heap of a fully
 commutative element in $W$ if and only if in the lattice embedding of
 $P$ in $S\times \N$,
 \begin{enumerate} 
 \item No column contains two points connected by an edge.
 \item For every pair $s,t\in S$ such that $m(s,t)\ge 3$, whenever there is a
 chain of edges connecting a sequence $s,t,s,\ldots $ of $m(s,t)$
 points, there is another chain connecting two points in this sequence.
 \end{enumerate} 
 \end{prop} 
 \noindent For example, from Figure~\ref{fig:example}, we easily see that
 the element $abcabd$ is fully commutative. In particular, although the
 heap contains the chains with labels $(a,b,a,b)$ and $(b,c,b)$ with $m(a,b)$ and $m(b,c)$ letters, respectively, each of these chains contains two points connected by the other chain. 
